\section{Coefficient decomposition and lifting for a ramified letter}

% BEGIN V81P1
\paragraph{V81P1. The integral deep and exceptional summands.}
Use an actual row of H80M2 and A81R4, with coefficient field $L$,
integers $O$, height $h=h'$, and letter $c$. Put $A=O[[v]]$,
$D(v)=E_c(v)$ from A81R3, $N_c=(e_cN(T))^\Delta$, and
$U=N_{c-1}/vN_{c-1}$. A81R7 identifies $U$ with the native
$L_c$ and makes it finite free over $O$. For each of the two
embeddings $\tau$ extending $c$, the map $U_L\to D_{\rm dR,\tau}$
taking a lifted vector to its native $\mathrm{valAt}$ and projecting
to $\tau$ is an isomorphism. In fact the preceding letter at
$X=0$ gives the standard lattice in crystalline coordinates by
W79H3. The primitive expansion of its Frobenius substitutes $pY$
in this expansion and applies the invertible crystalline Frobenius
from letter $c-1$ to $c$. Its constant coefficient is therefore
independent of the primitive root. This computes the native value:
multiplication by $Q^k$ shifts the leading coefficient by $k$ and
multiplies it by the $k$th power of the linear coefficient of $Q$,
which the definition divides out. Constants of an invariant
basis give an $L$-basis, by A81R7 and W79H3. Finally the crystalline
component projects isomorphically to each de Rham component by W79H1.
Pull the filtration back to $U_L$ by this map, writing $F_\tau^i$.

For each positive deep interval $[b_\tau,b_\tau+S]$, let
$C_\tau=F_\tau^{b_\tau}$. Its codimension is at most $m$, since
only the central jumps can lie below $b_\tau$. For each negative
deep interval $[-b_\tau-S,-b_\tau]$, put
$B_\tau=F_\tau^{1-b_\tau}$, of dimension at most $m$ by the
same central dimension bound. Let $B$ be the sum of the latter
subspaces and $C$ the intersection of the former. Split the
saturated submodule $U\cap B$ off $U$, and in its integral
complement $V_0$ split off $U_0=V_0\cap C$. These splittings
exist since the quotients are finite torsion-free over the DVR.
Write $U=U_0\oplus U_r$, placing $U\cap B$ in $U_r$.
The remaining summand of $U_r$ has rank at most
$\operatorname{codim} C$ by its map to $U_L/C$. Consequently
$\operatorname{rank} U_r\le2m$. Projection $P:U\to U_0$
kills every $B_\tau$ after inverting $p$, since $U\cap B$ spans
$B$. Lift a basis $e_j$ of $U_0$ to the preceding Wach letter.
Every such lift satisfies the native positive filtration conditions.
The construction allows either collection of signs to be empty.
% END V81P1

% BEGIN V81P2
\paragraph{V81P2. The comparison of the two integral models for a given frame.}
On this letter let $\Lambda= D(v)^{-h}A\otimes_O U$.
This is the source's $\mathrm{LamLA}$ on the full coefficient
lattice, with coefficient functions supported on $c$. At a
primitive root use $1+X=\zeta\exp(Y)$; then $v$ is a unit in
$L[[Y]]$ and $D$ is $Y$ times a unit. In fact $v-r_\tau$, where
$r_\tau=v(\zeta-1)$, is itself a uniformizer of this Taylor ring:
$D$ has a simple zero by $D\sim Q$ and the cyclotomic polynomial,
and its other linear factor is nonzero at this root.
In the constant coordinates of V81P1, the native submodule
$\mathcal N_L=\mathrm{NL}$ consists exactly of those $x\in\Lambda$
whose Laurent coefficient $x_i$ at each root lies in $F_\tau^{-i}$.
Indeed the defining numerator $D^hx$ tests its coefficient at $j$
against $F_\tau^{h-j}$. Multiplication by a unit in $L[[Y]]$ and
its inverse preserve this weighted condition, since earlier
coefficients lie in smaller filtered subspaces. The same reasoning
shows that $\mathcal N_L$ is stable under $A$ and $\gamma$.
On the physical letter write $P_c$ for the invariant Frobenius
span and $\Lambda_P=D^{-h}P_c$. For $y\in\Lambda_P$ we have
$D^{2h}y\in N_c$ by W79H4. Such a $y$ belongs to $N_c$ exactly
when all its Taylor coefficients belong to the corresponding
$F_\tau^{-i}$. W79H2--W79H4 give the filtered Laurent lattice
and detect the integral membership after clearing this power of
$Q\sim D$, on all primitive roots.

Fix any frame $(\alpha,\beta)$ as in the original propositions; P81F5 constructs one on these rows.
Its maps on the two quotients are $A$-linear, $\gamma$-equivariant,
natural for the row map, and their composites multiply by
$p^{2\sigma}$. The error modules on these quotients are
$J_L=v^M D^{3h}(A\otimes U)$ and $J_P=v^M D^{3h}P_c$.
They lie in $v^j\mathcal N_L$ and $v^jN_c$, respectively,
for $0\le j\le M$. On the left use the triple zero and
$F_\tau^{-h}=U_L$, and on the right use $D^hP_c\subset N_c$.
A81R7 also gives $N_c\subset\Lambda_P$.

Suppose $y$ represents $\alpha(x)$. The Taylor clause of the frame
says that the expansions of $y$ in period coordinates and of $p^\sigma x$
coincide in every degree $i<3h$. To see the cutoff, first multiply
by $Q^kD^h$, with $k$ clearing the physical vector as required
in the definition. The asserted numerator agreement holds in
degrees $<k+4h$. The multiplying series has leading degree $k+h$
with nonzero coefficient, so division gives the stated agreement.
For $i\ge3h$ the filtration test is automatic for both vectors.
The physical coefficients belong to the same crystalline component
by W79H2 after clearing the denominator. Thus $y\in N_c$ if and
only if $x\in\mathcal N_L$. For a representative $x$ of $\beta(y)$
use $\alpha(x)=[p^{2\sigma}y]$: the same Taylor test allows
division by this nonzero scalar in the filtered $L$-spaces.
This use of scalar division is only a rational filtration test
inside the respective integral modules where $x,y$ already lie.

The analogous implications for $v^j\mathcal N_L$ and $v^jN_c$
in their respective directions, $1\le j\le M$, follow by lifting
the unmultiplied representative, applying $A$-linearity and adding
the appropriate error module. Applying this at $j=1$ to
$\gamma-1$ proves that $\alpha$ and $\beta$ carry invariant
classes to invariant classes: explicitly these are the elements
of the respective $\Lambda$ for which the difference lies in
$v\mathcal N_L$ or $vN_c$. This establishes the needed
$\mathrm{ModelF}$ comparisons for any supplied frame.
% END V81P2

% BEGIN V81P3
\paragraph{V81P3. The first frame proposition in both positions.}
Let $w$ be an invariant physical vector at depth one, and let
$x\in\Lambda$ represent $\beta(w)$. By V81P2,
$(\gamma-1)x\in v\mathcal N_L$. At a primitive expansion
$\gamma$ acts as $Y\mapsto uY$, $u=\chi(\gamma)$.
Thus for $i\ne0$ we have $x_i\in F_\tau^{-i}$: the scalar
$u^i-1$ is nonzero, and multiplication by the Taylor unit $v$
preserves the weighted filtration. Let $\Omega$ be the union of
the integer intervals $[-b_\tau-S,-b_\tau)$ at positive labels
and $[b_\tau,b_\tau+S)$ at negative labels. They are contained
in $[-h,h]$ and exclude zero. Put
$H(Z)=\prod_{i\in\Omega}(Z-u^i)$, so $H(1)=\pi_{\Omega,1}$.

Apply the integral projection $P$ of V81P1 coefficientwise to $x$.
Write $f_j$ for the $e_j$-coordinate of $P H(\gamma)x$.
Each $D^hf_j$ is integral in $A$, as is
$D^h(\gamma-1)f_j/v$. Indeed $D^h x$ and
$D^h(\gamma-1)x/v$ are integral. Substitution by $\gamma$ takes
$v$ and $D$ to unit multiples of themselves, preserving both
integrality assertions under the polynomial and the projection.
The first unit assertion follows from the invariant coordinate of
A81R3; for the second, apply the substitution to $Q=\varphi(X)/X$
and use $\gamma(X)/X$ a unit and $D\sim Q$.

At a positive label, for $i<-b_\tau-S$ one has
$F_\tau^{-i}=0$. The remaining degrees $i<-b_\tau$ are killed
by $H(u^i)$. At a negative label, for $i<b_\tau$, $i\ne0$,
we have $x_i\in B_\tau$, which $P$ kills; the subsequent degrees
$i<b_\tau+S$ are killed by $H(u^i)$. Therefore the Taylor
expansion of $f_j$ starts in degree $-b_\tau$ at a positive
label, and at a negative label it agrees with a constant through
all degrees $<b_\tau+S$. Multiplication by $D^h$ gives exactly
both numerator tests of $\mathrm{IsCoord}$. The integral
difference established above supplies its last test.

Each $f_j\otimes e_j$ belongs to the unrestricted invariant module.
In fact the expansion of $(\gamma-1)f_j/v$ still starts in degree
$-b_\tau$ at a positive label, where $e_j\in F_\tau^{b_\tau}$.
At a negative label the constant term cancels before division, so
the expansion starts in degree $b_\tau+S$, and
$F_\tau^{-b_\tau-S}=U_L$. These give every required filtration
test. Write $x_r=(1-P)H(\gamma)x$. Its numerator belongs to
$A\otimes U_r$, and its difference lies in $v\mathcal N_L$ by
subtracting the deep summands from $H(\gamma)x$. Thus $x_r$ is
in the invariant module with coefficients restricted to $U_r$.
Finally $H(\gamma)x-H(1)x\in v\mathcal N_L$, since
$H(Z)-H(1)$ is divisible by $Z-1$ over $O$ and $v\mathcal N_L$
is stable under $\gamma$. Take $H(1)x$ as the representative
of $\beta(\pi_{\Omega,1}w)$. The preceding equality modulo
$v\mathcal N_L$, the lifts of $e_j$ and $\operatorname{rank}U_r
\le 2m$ give all clauses of $\mathrm{FirstPropositionFrame}$.
% END V81P3

% BEGIN V81P4
\paragraph{V81P4. The third slot proposition at either pure sign.}
Take the given onto row map $F:T_1\to T_2$ and a frame as above.
Let $f$ satisfy the native coordinate tests, $z\in N_{c-1}(T_2)$
satisfy the positive filtration tests, and write $\bar z$ for its
class at $v=0$. Take any representative $y$ of
$\alpha_2(f\otimes\bar z)$. We prove the original lifting
assertion with $\rho$ of A81R4.

When both intervals are positive, $f$ has Laurent order at least
$-b_\tau$ on each branch and $\bar z\in F_\tau^{b_\tau}$.
Thus $f\otimes\bar z\in\mathcal N_L$ and V81P2 puts $y$ in the
physical invariant Wach lattice $N_c(T_2)$.
By A81R9 choose $y_1\in N_c(T_1)$ with $Fy_1=p^\rho y$.
Both $y$ and $y_1$ are in their respective depth-one invariant
modules: membership in the physical denominator lattice and
triviality of $\gamma$ modulo $v$ are A81R7.

When both intervals are negative, write the expansion of $f$ at
each branch as a constant plus terms of degree at least $b_\tau+S$.
This follows by dividing the negative coordinate test by $D^h$,
whose order there is $h$. Since $\gamma$ kills the constant in
the difference, $(\gamma-1)f/v$ has the same lower bound on its
order. Its numerator after multiplication by $D^h$ is integral
by the final coordinate test. At either the source or the target
all vectors lie in $F_\tau^{-b_\tau-S}$ by the jump bound.
Therefore $f\otimes\bar z'$ belongs to the invariant coefficient
module for every $z'$ of the preceding invariant Wach letter,
in either representation.

Now use A81R9 on the preceding letter to lift $p^\rho z$ to
$z_1$. Choose $y_1$ representing $\alpha_1(f\otimes\bar z_1)$.
V81P2 puts $y_1$ and $y$ in the respective depth-one physical
invariant modules. Naturality of the given frame and its linearity
give $F y_1-p^\rho y\in J_P\subset vN_c(T_2)$.
This is the native congruence in $\mathrm{LiftsBy}$, for every
representative $y$. Together the two cases give the pure-sign
$\mathrm{ThirdPropositionSlot}$ on the ramified rows at $q$ with
its prescribed exponent.
% END V81P4
